\section*{附录 A: 配置文件清单}

本项目采用 YAML 格式进行配置管理，以下列出所有配置文件及其详细说明。

\subsection*{主实验配置}

所有主实验配置均位于 \texttt{config/} 目录下。

\subsubsection*{exp\_rna.yaml（RNA 单模态实验）}

\begin{verbatim}
experiment_name: "RNA Only"
data:
  use_rna: true
  use_meth: false
  use_feature_selection: true
rna_preprocessing:
  log_transform: true
  mean_filter_threshold: 0.5
  variance_cutoff: 0.95
svm:
  kernel: "rbf"
  C: 1.0
  random_state: 42
\end{verbatim}

该配置专用于基于 RNA-seq 数据的基准实验。特征选择采用方差排序，保留前 95\% 的方差特征。SVM 采用 RBF 核，正则化参数 $C=1.0$。

\subsubsection*{exp\_concat.yaml（早期融合实验）}

\begin{verbatim}
experiment_name: "Concatenation"
data:
  use_rna: true
  use_meth: true
  use_feature_selection: true
rna_preprocessing:
  log_transform: true
  mean_filter_threshold: 0.5
  variance_cutoff: 0.95
meth_preprocessing:
  variance_cutoff: 0.95
svm:
  kernel: "rbf"
  C: 1.0
  random_state: 42
\end{verbatim}

该配置将 RNA 和甲基化数据在预处理后进行直接拼接。两个模态均按方差排序进行特征选择，确保输入维度一致性。

\subsubsection*{exp\_mofa.yaml（晚期融合实验）}

\begin{verbatim}
experiment_name: "MOFA"
data:
  use_rna: true
  use_meth: true
  use_feature_selection: true
rna_preprocessing:
  log_transform: true
  mean_filter_threshold: 0.5
  variance_cutoff: 0.95
meth_preprocessing:
  variance_cutoff: 0.95
mofa:
  n_factors: 10
  max_iteration: 100
  convergence_mode: "slow"
svm:
  kernel: "rbf"
  C: 1.0
  random_state: 42
\end{verbatim}

该配置采用 MOFA 进行潜在因子分解，提取 10 个潜在因子用于后续 SVM 分类。MOFA 的训练采用缓慢收敛模式以确保稳定性。

\subsection*{消融实验配置}

\subsubsection*{ablation\_no\_meth.yaml（去甲基化）}

\begin{verbatim}
experiment_name: "No Methylation"
data:
  use_rna: true
  use_meth: false
  use_feature_selection: true
\end{verbatim}

该配置移除甲基化数据，仅使用 RNA 信息进行融合。用于评估甲基化对模型性能的贡献。

\subsubsection*{ablation\_no\_rna.yaml（去 RNA）}

\begin{verbatim}
experiment_name: "No RNA"
data:
  use_rna: false
  use_meth: true
  use_feature_selection: true
\end{verbatim}

该配置移除 RNA 数据，仅使用甲基化信息。用于评估 RNA 对融合模型的贡献度。

\subsubsection*{ablation\_no\_fs.yaml（无特征选择）}

\begin{verbatim}
experiment_name: "No Feature Selection"
data:
  use_rna: true
  use_meth: true
  use_feature_selection: false
\end{verbatim}

该配置禁用特征选择，使用所有原始特征进行融合。用于验证特征选择的必要性。

\section*{附录 B: 项目结构}

\begin{verbatim}
Cancer-Classification/
├── config/                    # 实验配置文件
│   ├── exp_rna.yaml
│   ├── exp_concat.yaml
│   ├── exp_mofa.yaml
│   ├── ablation_no_meth.yaml
│   ├── ablation_no_rna.yaml
│   └── ablation_no_fs.yaml
├── datasets/                  # 数据文件
│   ├── TCGA-BRCA.star_counts.csv
│   ├── TCGA-BRCA.methylation450.csv
│   └── brca_pam50.csv
├── src/                       # 源代码
│   ├── pipeline.py           # 主流程
│   ├── data/
│   │   ├── loader.py         # 数据加载
│   │   ├── align.py          # 样本对齐
│   │   └── preprocess.py     # 数据预处理
│   ├── features/
│   │   └── mofa.py           # MOFA 包装器
│   ├── models/
│   │   └── train.py          # 模型训练
│   └── visualization/        # 可视化模块
│       ├── tsne.py
│       ├── feature_distribution.py
│       ├── model_performance.py
│       ├── fusion_methods.py
│       └── generate_report_figures.py   # 兼容包装器，现已转向白名单图表流水线
├── experiments/
│   └── run.py                # 实验运行脚本
├── outputs/
│   ├── figures/              # 生成的图表
│   └── logs/                 # 实验日志
├── run_all.sh                # 主实验批处理脚本
├── scripts/
│   └── run_ablation.sh       # 消融实验脚本
└── docs/
    └── 研究报告/             # 论文源文件
\end{verbatim}

\section*{附录 C: 数据说明}

\subsection*{TCGA-BRCA 数据集}

TCGA-BRCA 数据集包含 1,100+ 乳腺癌患者的多组学数据，本研究使用的子集包括：

\begin{itemize}
    \item \textbf{RNA-seq 数据}：20,531 个基因的表达量，来自 200 个样本
    \item \textbf{DNA 甲基化数据}：450,000 个 CpG 位点的甲基化水平，来自 200 个样本
    \item \textbf{临床标签}：基于 PAM50 分类的四个乳腺癌亚型
    \begin{itemize}
        \item Luminal A：45 个样本
        \item Luminal B：52 个样本
        \item HER2+：38 个样本
        \item Basal：45 个样本
    \end{itemize}
\end{itemize}

\subsection*{数据预处理流程}

\begin{enumerate}
    \item \textbf{样本对齐}：利用 RNA 和甲基化数据的样本交集，进行确定性对齐
    \item \textbf{数据清洗}：移除重复样本，删除缺失值
    \item \textbf{特征选择}：基于方差排序，RNA 保留前 95\% 方差特征，甲基化同样
    \item \textbf{标准化}：采用 StandardScaler 对所有特征进行零均值单位方差标准化
    \item \textbf{列车测试分割}：80\% 训练，20\% 测试，分割后再进行预处理以避免泄露
\end{enumerate}

\section*{附录 D: 运行指南}

\subsection*{环境配置}

\begin{verbatim}
# 安装依赖（已在 requirements.txt 中列出）
pip install pandas numpy scikit-learn mofapy2 matplotlib seaborn

# 或使用项目自带的 Python 环境
./medical/bin/python3 -m pip install -r requirements.txt
\end{verbatim}

\subsection*{运行主实验}

\begin{verbatim}
# 方式 1：使用批处理脚本
bash run_all.sh

# 方式 2：单独运行某个实验
./medical/bin/python3 experiments/run.py --config config/exp_rna.yaml

# 方式 3：使用 Python 直接调用
./medical/bin/python3 -c "from src.pipeline import run_pipeline; \
    run_pipeline('config/exp_rna.yaml')"
\end{verbatim}

\subsection*{生成可视化图表}

\begin{verbatim}
# 生成论文白名单图表并归档非正文图
./medical/bin/python3 scripts/generate_paper_artifacts.py --archive-legacy

# 输出目录：outputs/figures/
# 包含：统计主结果图、稳定性图、SHAP summary
\end{verbatim}

\subsection*{预期输出}

所有实验的结果均输出到 \texttt{outputs/logs/} 目录，格式如下：

\begin{verbatim}
Experiment: RNA Only
Accuracy: 0.750
Precision: 0.760
Recall: 0.740
F1-Score: 0.750

Experiment: Concat
Accuracy: 0.875
Precision: 0.880
Recall: 0.870
F1-Score: 0.875

...
\end{verbatim}

\section*{附录 E: 扩展功能}

\subsection*{自定义特征选择}

用户可以在 \texttt{src/data/preprocess.py} 中修改 \texttt{variance\_cutoff} 参数，调整特征选择阈值：

\begin{verbatim}
# 保留前 90\% 方差特征
variance_cutoff: 0.90

# 保留前 98\% 方差特征
variance_cutoff: 0.98
\end{verbatim}

\subsection*{SVM 超参数调优}

在 YAML 配置中修改 SVM 参数：

\begin{verbatim}
svm:
  kernel: "rbf"      # 可选: linear, poly, rbf, sigmoid
  C: 1.0             # 正则化参数，取值 (0, inf)
  gamma: "scale"     # RBF 核参数
  degree: 3          # poly 核的多项式度数
  random_state: 42
\end{verbatim}

\subsection*{MOFA 参数配置}

调整 MOFA 潜在因子数量和收敛参数：

\begin{verbatim}
mofa:
  n_factors: 15              # 潜在因子数量
  max_iteration: 200         # 最大迭代次数
  convergence_mode: "fast"   # 收敛模式：fast, medium, slow
  seed: 42
\end{verbatim}

\section*{附录 F: 故障排除}

\subsection*{常见问题与解决方案}

\begin{enumerate}
    \item \textbf{问题}：\texttt{KeyError: 'sample'}
    
    \textbf{原因}：临床数据文件格式不匹配，预期 TSV 格式，实际为其他格式
    
    \textbf{解决}：检查 \texttt{datasets/brca\_pam50.csv} 的分隔符，确保为 Tab 分隔
    
    \item \textbf{问题}：MOFA 训练失败，\texttt{IndexError: tuple index out of range}
    
    \textbf{原因}：MOFA 输入数据格式不正确，或样本数过少
    
    \textbf{解决}：确保在 \texttt{src/features/mofa.py} 中数据格式为 nested list（views × groups）
    
    \item \textbf{问题}：内存不足，处理大规模甲基化数据
    
    \textbf{原因}：甲基化数据维度过高（450,000 特征）
    
    \textbf{解决}：增加 \texttt{variance\_cutoff} 参数值，保留更少特征；或采用样本子集进行实验
    
    \item \textbf{问题}：模型精度低于预期
    
    \textbf{原因}：可能是特征选择不当，或 SVM 超参数未调优
    
    \textbf{解决}：进行网格搜索或随机搜索调优 SVM 参数 $C$ 和 $\gamma$
\end{enumerate}

\subsection*{性能优化建议}

\begin{itemize}
    \item 对于大规模数据，考虑使用 GPU 加速的库（如 CuPy、RAPIDS）
    \item 使用并行处理加速特征选择：\texttt{n\_jobs=-1}
    \item 对甲基化数据进行增量处理，避免一次性加载所有特征
    \item 利用 joblib 进行模型训练的并行化
\end{itemize}
